\documentclass[12pt]{book}
\usepackage{precalc1}
\setcounter{chapter}{4}
\begin{document}

\chapter{Inverses and Radicals}

A function sends inputs forward to outputs. This week we ask when that motion can be
reversed. If an output remembers exactly which input produced it, then the function can
be undone. If two different inputs produce the same output, the reverse direction has a
choice to make, and that choice prevents the reverse relation from being a function.

This idea connects several topics that may at first look separate. Domain and range tell
us where a function starts and where it lands. One-to-one functions are the functions
whose outputs do not get shared. Inverses run one-to-one functions backward. Radicals,
such as square roots and cube roots, are inverses of power functions, so their domains
come from the same reversing process. The goal is to see these as one story rather than
as four unrelated rules.

\section{Domain and Range}

Before a function can be reversed, we need to know its input set and output set. The
\emph{domain} is the set of inputs the function is allowed to accept. The \emph{range}
is the set of outputs the function actually produces. These two sets are part of the
function's meaning. A formula by itself tells us how to compute, but the domain and
range tell us where the computation is allowed to live.

For a function given by a formula, the natural domain is every real input for which the
formula is defined. The range is usually read from the graph, from transformations, or
from the nature of the formula. In this course the most common domain restrictions come
from two operations: division and even roots.

\begin{definition}{Domain and range}
The \emph{domain} of $f$ is the set of all inputs $x$ for which $f(x)$ is defined.
The \emph{range} is the set of all outputs $f(x)$ produced as $x$ ranges over the
domain.
\end{definition}

Division forbids a zero denominator, because division by zero is not defined. An even
root forbids a negative radicand, because the even root of a negative number is not real.
Everything else we use here, such as polynomials, odd roots, and absolute values, accepts
all real inputs unless an outside restriction is stated.

\begin{keyidea}
To find a natural domain from a formula, look for forbidden operations:
\begin{itemize}\setlength{\itemsep}{2pt}
  \item denominators cannot equal zero;
  \item radicands under even roots must be nonnegative;
  \item odd roots, polynomials, and absolute values do not create domain restrictions
        by themselves.
\end{itemize}
The range then records the outputs that actually occur after the domain is chosen.
\end{keyidea}

\begin{example}[reading the domain from the formula]
For $f(x)=\dfrac{1}{x-5}$ the denominator vanishes at $x=5$, so the domain is all
reals except $5$. For $g(x)=\sqrt{x-3}$ the radicand $x-3$ must be nonnegative, so the
domain is $x\ge 3$. For $h(x)=x^2+1$ no operation restricts the input, so the domain is
all reals, while the range is $y\ge 1$ because a square is never negative.
\end{example}

The important point in the example is that the domain is not guessed from the general
shape. It is forced by the operations in the formula. The range asks a different
question: after the allowed inputs are used, which outputs appear?

\begin{pitfall}
Domain restrictions come from the operations, not from the graph's appearance. An odd
root such as $\sqrt[3]{x}$ accepts negatives, since a cube root of a negative number is
a real negative number; only \emph{even} roots restrict the radicand.
\end{pitfall}

\section{One-to-One and Onto}

Now we turn from the question ``where is the function defined?'' to the question ``can
its arrows be reversed?'' A function already has the rule that each input goes to exactly
one output. For an inverse function, we need the reverse assignment also to be single
valued: each output must point back to exactly one input.

Two properties describe how a function distributes its inputs over its outputs. The first,
one-to-one, says that no two inputs share an output. The second, onto, says that every
output in a chosen target set is hit by the function.

\begin{definition}{One-to-one and onto}
A function $f$ is \emph{one-to-one} if distinct inputs always give distinct outputs:
$x_1\neq x_2$ forces $f(x_1)\neq f(x_2)$. It is \emph{onto} a target set of outputs
if every value in that set is attained by some input.
\end{definition}

One-to-one is the property that matters most for inverses in this course. If an output is
shared, the reverse direction cannot decide which input to choose. The graph gives a
clean test, because a horizontal line collects all points with the same output value.

\begin{keyidea}
\textbf{Horizontal line test.} A function is one-to-one if and only if no horizontal
line crosses its graph more than once. A repeated crossing marks two inputs sharing
one output, which violates the one-to-one property.
\end{keyidea}

\begin{center}
\begin{tikzpicture}[scale=0.8]
  % one-to-one: a cubic
  \begin{scope}
    \draw[->] (-2,0)--(2,0) node[right]{\scriptsize $x$}; \draw[->] (0,-2.2)--(0,2.2);
    \draw[pcblue,thick,domain=-1.5:1.5,smooth,samples=60] plot (\x,{0.6*\x*\x*\x});
    \draw[pcred,dashed] (-1.9,0.9)--(1.9,0.9);
    \fill[pcred] (1.106,0.9) circle (1.6pt);
    \node at (0,-3) {\scriptsize one-to-one};
  \end{scope}
  % not one-to-one: a parabola
  \begin{scope}[xshift=6cm]
    \draw[->] (-2,0)--(2,0) node[right]{\scriptsize $x$}; \draw[->] (0,-0.4)--(0,3.4);
    \draw[pcblue,thick,domain=-1.7:1.7,smooth] plot (\x,{\x*\x});
    \draw[pcred,dashed] (-1.9,1.4)--(1.9,1.4);
    \fill[pcred] (1.183,1.4) circle (1.6pt);
    \fill[pcred] (-1.183,1.4) circle (1.6pt);
    \node at (0,-1) {\scriptsize not one-to-one};
  \end{scope}
\end{tikzpicture}
\end{center}

\begin{example}[applying the horizontal line test]
The line $f(x)=2x-1$ is one-to-one; every horizontal line meets it once. The parabola
$f(x)=x^2$ is not one-to-one; the horizontal line $y=4$ meets it at $x=2$ and
$x=-2$, two inputs sharing the output $4$. Restricting the parabola to $x\ge0$ removes
the left half and makes the restricted function one-to-one.
\end{example}

The restriction in the last sentence is not a trick. It is how many inverses are created:
if the original function repeats outputs, we sometimes choose a domain where the repeats
are removed. The square root function comes from exactly this choice.

Whether a function is onto depends on the target set declared for its outputs. The
squaring function is onto the nonnegative reals (every nonnegative number is a square)
but not onto all reals (no input squares to a negative number). Choosing the target set
to be exactly the range makes any function onto by construction.

\newpage
\section{The Inverse Function}

When $f$ is one-to-one, each output came from exactly one input, so the assignment can be
reversed: send each output back to the input that produced it. That reverse assignment
is the inverse function $f^{-1}$. The notation looks like an exponent, but here it means
``inverse function,'' not reciprocal.

\begin{definition}{Inverse function}
If $f$ is one-to-one, its \emph{inverse} $f^{-1}$ is the function satisfying
\[
f^{-1}(y)=x \quad\text{exactly when}\quad f(x)=y.
\]
Equivalently, $f^{-1}(f(x))=x$ for every $x$ in the domain of $f$, and
$f(f^{-1}(y))=y$ for every $y$ in the range.
\end{definition}

\begin{keyidea}
An inverse undoes the function. Its domain is the range of $f$ and its range is the
domain of $f$; the two sets trade places, because reversing the arrows swaps inputs
and outputs. Only a one-to-one function has an inverse, since a shared output would
have no single input to return to.
\end{keyidea}

\subsection{Finding an inverse}

Because the inverse swaps the roles of input and output, the algebra is: write
$y=f(x)$, exchange $x$ and $y$, and solve for $y$. This procedure is not just a formal
recipe. The swap represents the reversal of the arrows. Solving for $y$ puts the
reversed rule back into function notation.

\begin{example}[finding an inverse]
Invert $f(x)=2x-1$. Write $y=2x-1$, swap to $x=2y-1$, and solve:
\[
y=\frac{x+1}{2}.
\]
So $f^{-1}(x)=\dfrac{x+1}{2}$. Check:
\[
f^{-1}(f(x))=\frac{(2x-1)+1}{2}=x.
\]
The check confirms that applying $f$ and then $f^{-1}$ brings the input back to where
it started.
\end{example}

\begin{visual}
Swapping $x$ and $y$ reflects the graph across the line $y=x$. The graph of $f^{-1}$
is the mirror image of the graph of $f$ in that diagonal, because every point $(a,b)$
on $f$ becomes the point $(b,a)$ on $f^{-1}$.
\end{visual}

\begin{center}
\begin{tikzpicture}[scale=0.75]
  \draw[->] (-0.5,0)--(5,0) node[right]{\scriptsize $x$};
  \draw[->] (0,-0.5)--(0,5) node[above]{\scriptsize $y$};
  \draw[gray,dashed] (-0.4,-0.4)--(4.8,4.8) node[right]{\scriptsize $y=x$};
  \draw[pcblue,thick,domain=0:2.1,smooth] plot (\x,{\x*\x});
  \node[pcblue] at (1.7,4.4) {\scriptsize $f(x)=x^2$};
  \draw[pcred,thick,domain=0:4.4,smooth] plot (\x,{sqrt(\x)});
  \node[pcred] at (4.4,1.6) {\scriptsize $f^{-1}(x)=\sqrt{x}$};
\end{tikzpicture}
\end{center}

The picture also shows why domain restrictions matter. The full parabola $y=x^2$ would
reflect to a sideways parabola, which fails the vertical line test. The restricted right
half reflects to the square-root graph, which is a function.

\section{Radicals as Inverses of Powers}

A radical is the inverse of a power function. The square root inverts squaring, the cube
root inverts cubing, and in general $\sqrt[n]{x}$ inverts $x^n$. The parity of $n$
decides whether the power function repeats outputs. Odd powers keep moving in one
direction, so they are one-to-one on all real numbers. Even powers fold the line across
the $y$-axis, sending $x$ and $-x$ to the same output, so they must be restricted before
being inverted.

\begin{keyidea}
\textbf{Radicals and parity.}
\begin{itemize}\setlength{\itemsep}{2pt}
  \item \emph{$n$ odd.} $x^n$ is one-to-one on all reals, so $\sqrt[n]{x}$ is defined
        for all reals; negatives are allowed.
  \item \emph{$n$ even.} $x^n$ is not one-to-one on all reals (it shares outputs
        across $\pm x$), so it must be restricted to $x\ge0$ before inverting.
        Its inverse $\sqrt[n]{x}$ then has domain and range $x\ge0$.
\end{itemize}
\end{keyidea}

\begin{example}[why the square root needs a restriction]
The function $x^2$ fails the horizontal line test on all reals. Restricted to $x\ge0$
it is one-to-one, and its inverse is $\sqrt{x}$, with domain $x\ge0$ and range
$y\ge0$. The discarded left half $x\le0$ would have inverted to $-\sqrt{x}$; the two
halves are the two branches of the sideways parabola, only one of which is a function.
\end{example}

\begin{example}[an odd root accepts negatives]
The function $x^3$ is one-to-one on all reals, so $\sqrt[3]{x}$ is defined everywhere.
For instance $\sqrt[3]{-8}=-2$, since $(-2)^3=-8$. No restriction is needed because
cubing never repeats an output.
\end{example}

\begin{center}
\begin{tikzpicture}[scale=0.7]
  % cube root, all reals
  \begin{scope}
    \draw[->] (-2.4,0)--(2.4,0) node[right]{\scriptsize $x$};
    \draw[->] (0,-1.8)--(0,1.8) node[above]{\scriptsize $y$};
    \draw[pcblue,thick,domain=-1.25:1.25,smooth,samples=80,variable=\t]
      plot ({\t*\t*\t},{\t});
    \node at (0,-2.6) {\scriptsize $\sqrt[3]{x}$, all reals};
  \end{scope}
  % square root, x>=0
  \begin{scope}[xshift=7cm]
    \draw[->] (-0.6,0)--(4,0) node[right]{\scriptsize $x$};
    \draw[->] (0,-0.6)--(0,2.4) node[above]{\scriptsize $y$};
    \draw[pcred,thick,domain=0:3.6,smooth,samples=60] plot (\x,{sqrt(\x)});
    \node at (1.8,-1.2) {\scriptsize $\sqrt{x}$, $x\ge0$};
  \end{scope}
\end{tikzpicture}
\end{center}

\subsection{Solving equations with radicals}

A radical equation is solved by isolating the radical and raising both sides to the
matching power. The solving step is allowed, but it is not always reversible. Raising an
even root to an even power can introduce a false solution, because the squaring step
forgets whether the original right-hand side was nonnegative. That is why every
candidate must be checked in the original equation.

\begin{example}[solving and checking]
Solve $\sqrt{x+5}=x-1$. Square both sides:
\[
x+5=x^2-2x+1,
\]
so $x^2-3x-4=0$, giving $(x-4)(x+1)=0$ and candidates $x=4$ and $x=-1$. Check in the
original: $x=4$ gives $\sqrt{9}=3$ and $4-1=3$, valid. $x=-1$ gives $\sqrt{4}=2$ but
$-1-1=-2$; the two sides disagree, so $x=-1$ is rejected. The only solution is
$x=4$.
\end{example}

\begin{pitfall}
Squaring both sides can manufacture solutions the original equation never had,
because $\sqrt{\,}$ returns the nonnegative root while squaring forgets the sign.
Always substitute each candidate back into the original radical equation and discard
any that fail.
\end{pitfall}

\section{Exercises}

\begin{exercises}
\item State the domain of each: (a) $\dfrac{x}{x^2-9}$; (b) $\sqrt{2x-6}$;
(c) $\sqrt[3]{x+1}$; (d) $\dfrac{1}{\sqrt{x-4}}$.

\item Give the range of each, with a reason: (a) $f(x)=x^2-4$; (b) $f(x)=\sqrt{x}+2$;
(c) $f(x)=-|x|$.

\item Apply the horizontal line test to decide which are one-to-one:
(a) $f(x)=3x+2$; (b) $f(x)=x^2-1$; (c) $f(x)=x^3$; (d) $f(x)=|x|$.

\item Explain, using the definition, why a function must be one-to-one to have an
inverse. What goes wrong at a shared output?

\item Find the inverse of each and state its domain: (a) $f(x)=3x-7$;
(b) $f(x)=\dfrac{x+2}{5}$; (c) $f(x)=x^3+1$.

\item For $f(x)=x^2$ restricted to $x\ge0$, find $f^{-1}$, verify
$f^{-1}(f(x))=x$ on the restricted domain, and state the domain and range of $f^{-1}$.

\item Explain why $\sqrt[3]{x}$ is defined for all real numbers while $\sqrt{x}$ is
not, in terms of the parity of the power being inverted.

\item Sketch $f(x)=x^2$ on $x\ge0$ and its inverse on the same axes, and show that
each is the reflection of the other across $y=x$.

\item Solve and check, discarding any false solutions: (a) $\sqrt{2x+3}=x$;
(b) $\sqrt{x+7}=x+1$.

\item For an increasing one-to-one function, intersections of the graph of $f$ with
the graph of $f^{-1}$ occur on the line $y=x$. Explain this using the reflection
picture, and contrast it with a decreasing self-inverse function such as $f(x)=1-x$.

\item The function $f(x)=\sqrt{x-2}+1$ is built from the parent $\sqrt{x}$ by
transformations. State the domain and range, find the inverse, and confirm the domain
of the inverse equals the range of $f$.

\item Determine whether $f(x)=\dfrac{1}{x}$ is its own inverse by computing
$f^{-1}$, and explain the symmetry of its graph in light of the result.
\end{exercises}

\end{document}
